%%%%%%%%%%  scan idx 547  |  folio 534  %%%%%%%%%%
\SGAPrintPage{534}%

\[
\alpha:P\xrightarrow{i}A\xrightarrow{j}A\otimes_{\mathbb Z}\mathbb Q
\]
\uline{(where} $j(x)=x\otimes1$\uline{) is a polynomial map. If} $I^{m+1}=0$\uline{, then} $\alpha$ \uline{has degree} $\leq m$.

\emph{Proof}: Let $x_1,\ldots,x_r$ be a sequence of elements of $P$ with $x_i=1+y_i$ ($i=1,\ldots,r$). Then the binomial formula gives, for every $n_1,\ldots,n_r\in Z$:
\sgaStableTarget{sga6.e10.eq.2-1-5}\begin{equation}\tag{2.1.5}
x_1^{n_1}\cdots x_r^{n_r}=
\sum_{k_1,\ldots,k_r\geq0}\binom{n_1}{k_1}\cdots\binom{n_r}{k_r}y_1^{k_1}\cdots y_r^{k_r},
\end{equation}
where the sum on the right is finite because the elements $y_1,\ldots,y_r$ are nilpotent. But the coefficients $\binom{n_1}{k_1}\cdots\binom{n_r}{k_r}$ are polynomials with rational coefficients in the variables $n_1,\ldots,n_r$, of degree $k_1+\cdots+k_r$, so $\alpha$ is indeed a polynomial map. If $I^{m+1}=0$, then $y_1^{k_1}\cdots y_r^{k_r}=0$ for $k_1+\cdots+k_r>m$. Thus $\alpha$ has degree $\leq m$, Q.E.D.

\sgaStableTarget{sga6.e10.sec.2.2}\subsection*{2.2.}
Now consider a noetherian scheme $X$ and an element $y$ of $\Filt_j(X)$ (1.1.1). Associating to every $x\in K^{\bullet}(X)$ the element $xy$ of $K_\bullet(X)$ gives a homomorphism of $K^{\bullet}(X)$-modules
\[
\Theta_y:K^{\bullet}(X)\longrightarrow K_\bullet(X)
\]
which, taking account of (1.3.2), factors through the canonical epimorphism
\[
j:K^{\bullet}(X)\longrightarrow K^{\bullet}(X)/\Filt^{j+1}(X)=A.
\]


%%%%%%%%%%  scan idx 548  |  folio 535  %%%%%%%%%%
\SGAPrintPage{535}%

Let $\beta_y:A\to K_\bullet(X)$ be the homomorphism so defined. Denoting by $I$ the ideal $\Filt^1(X)/\Filt^{j+1}(X)$ of $A$, one has $I^{j+1}=0$. Keeping the notation of Lemma 2.1.4, one obtains a homomorphism of abelian groups
\[
\gamma:\Pic(X)\longrightarrow P
\]
that associates to the class in $\Pic(X)$ of every invertible $\underline{O}_X$-module $\underline{L}$ the element $j(\cl^{\bullet}(\underline{L}))$ of $P$.

Lemmas 2.1.2 and 2.1.4 immediately give:

\sgaStableTarget{sga6.e10.prop.2-2-1}\textbf{Proposition 2.2.1.} \uline{For every noetherian scheme} $X$ \uline{and every} $y\in\Filt_j(X)$ ($j\in\mathbb Z$\uline{), the composite map}
\[
\Pic(X)\longrightarrow K_\bullet(X)\otimes_{\mathbb Z}\mathbb Q,
\]
\uline{defined by}
\[
\cl_P(\underline{L})\rightsquigarrow (\cl^{\bullet}(\underline{L})y)\otimes1
\]
\uline{is a polynomial map of degree} $\leq j$.

\sgaStableTarget{sga6.e10.sec.2.3}\subsection*{2.3.}
If $X$ is a proper scheme over a field $k$ and $g:X\to\Spec(k)$ is the structural morphism, then the homomorphism
\[
g_*:K_\bullet(X)\longrightarrow K_\bullet(\Spec(k))\cong\mathbb Z
\]
associates to the class $\cl_\bullet(\underline{F})$ of every coherent $\underline{O}_X$-module $\underline{F}$ the Euler--Poincaré characteristic $\chi(\underline{F})\in\mathbb Z$ of $\underline{F}$ on $X$. From 2.2.1 one immediately concludes:


%%%%%%%%%%  scan idx 549  |  folio 536  %%%%%%%%%%
\SGAPrintPage{536}%

\sgaStableTarget{sga6.e10.prop.2-3-1}\textbf{Proposition 2.3.1.} \uline{For every proper scheme} $X$ \uline{over} $k$\uline{, and every} $y\in\Filt_j(X)$\uline{, the map}
\[
s_y:\Pic(X)\longrightarrow\mathbb Z\hookrightarrow\mathbb Q
\]
\uline{defined by}
\[
\cl_P(\underline{L})\rightsquigarrow g_*(\cl^{\bullet}(\underline{L})y)=\chi(\cl^{\bullet}(\underline{L})y)
\]
\uline{is a polynomial map of degree} $\leq j$.

When $\underline{L}_1,\ldots,\underline{L}_n$ are invertible $\underline{O}_X$-modules and $\underline{F}$ is a coherent $\underline{O}_X$-module such that $\dim(\Supp(\underline{F}))=j$, then 2.3.1. implies in particular that the Euler--Poincaré characteristic
\sgaStableTarget{sga6.e10.eq.2-3-2}\begin{equation}\tag{2.3.2}
\chi(\underline{L}_1^{\otimes m_1}\otimes\cdots\otimes \underline{L}_n^{\otimes m_n}\otimes \underline{F})
\end{equation}
is a numerical polynomial in the variables $m_1,\ldots,m_n\in\mathbb Z$, of degree $\leq\dimdot\operatorname{supp.}F$, sometimes called the \emph{Snapper polynomial} for $\underline{F}$ and the sequence $\underline{L}_1,\ldots,\underline{L}_n$.

\sgaStableTarget{sga6.e10.sec.2.4}\subsection*{2.4.}
Recall [2] that the Picard functor $\text{\uwave{Pic}}_{X/k}$ of $X$ over $k$ is the sheaf of groups for the ($\fpqc$) topology over $k$ associated with the presheaf
\[
T\rightsquigarrow \Pic(X\times_kT)\qquad (T\in\operatorname{Ob}(\Sch/k)).
\]
When $k$ is algebraically closed, the canonical homomorphism
\[
\Pic(X)\longrightarrow\Pic_{X/k}(\Spec(k))
\]
is known to be an isomorphism.

% Local continuation macros
\providecommand{\Filt}{\operatorname{Fil}}
\providecommand{\fpqc}{\mathrm{fpqc}}
\providecommand{\Sch}{\operatorname{Sch}}
\renewcommand{\Ob}{\operatorname{ob}}
\providecommand{\Gm}{\mathbb G_m}
\providecommand{\Der}{\operatorname{Der}}
\providecommand{\SL}{\operatorname{SL}}
\providecommand{\Aut}{\operatorname{Aut}}
\providecommand{\Mod}{\operatorname{Mod}}
\providecommand{\Tors}{\operatorname{Tors}}
\providecommand{\Mat}{\operatorname{Mat}}
\providecommand{\len}{\operatorname{long}}
\renewcommand{\id}{\operatorname{id}}
\providecommand{\Spec}{\operatorname{Spec}}
\providecommand{\Picf}{\operatorname{Pic}}
\providecommand{\Ens}{\operatorname{Ens}}
\providecommand{\clp}{\operatorname{cl}_{P}}
\providecommand{\clo}{\operatorname{cl}^{\circ}}
\providecommand{\clb}{\operatorname{cl}_{\bullet}}
\providecommand{\Grb}{\operatorname{Gr}_{\bullet}}
\providecommand{\Gro}{\operatorname{Gr}^{\bullet}}
\providecommand{\longO}[1]
{\operatorname{length}_{#1}}

\renewcommand{\Gro}{\operatorname{Gr}^{\bullet}}

%%%%%%%%%%  scan idx 550  |  folio 537  %%%%%%%%%%
\SGAPrintPage{537}%

When \(X\) is proper over \(k\), it is known [4] that \(\text{\uwave{Pic}}_{X/k}\) is representable by a group scheme \(\underline{\Pic}_{X/k}\) locally of finite type over \(k\). The neutral component \(\underline{\Pic}^0_{X/k}\) of \(\underline{\Pic}_{X/k}\) is a connected commutative algebraic group over \(k\).

\sgaStableTarget{sga6.e10.def.2-4-1}\textbf{Definition 2.4.1.} \uline{Let} \(X\) \uline{be a proper scheme over a field} \(k\), \uline{and} \(\underline{L}\) \uline{an invertible} \(\underline{O}_X\)\uline{-module} (resp. \(x\) \uline{an element of} \(\Pic(X)\)\uline{). We say that} \(\underline{L}\) (resp. \(x\)\uline{) is} \(\tau\)\uline{-equivalent to zero when there is an integer} \(n\ne0\) \uline{such that} \(\underline{L}^{\otimes n}\) (resp. \(nx\)\uline{) is algebraically equivalent to zero. The elements of} \(\Pic(X)\) \uline{algebraically equivalent} (resp. \(\tau\)\uline{-equivalent) to zero form a subgroup} \(\Pic^0(X)\) (resp. \(\Pic^\tau(X)\)\uline{) of} \(\Pic(X)\).

\sgaStableTarget{sga6.e10.lem.2-4-3}\textbf{Lemma 2.4.3.} \uline{For every proper scheme} \(X\) \uline{over an algebraically closed field} \(k\), \uline{the group} \(\Pic^0(X)\) \uline{is divisible by every integer} \(n\) \uline{prime to the characteristic of} \(k\).

\emph{Proof}: Since \(k\) is algebraically closed, one has \(\Pic^0(X)\cong \underline{\Pic}^0_{X/k}(k)\). Since \(\underline{\Pic}^0_{X/k}\) is connected and of finite type, the \(n\)-th-power endomorphism is surjective when \(n\) is prime to the characteristic of \(k\) (SGA 3 XV 1.3 and VI 1.3.2). But \(k\) is algebraically closed, so the corresponding endomorphism of the group \(\underline{\Pic}^0_{X/k}(k)\) is surjective, Q.E.D.

From 2.4.3 one concludes that, in the exact sequence
\[
0\longrightarrow\Pic^0(X)\longrightarrow\Pic^\tau(X)\longrightarrow T\longrightarrow0,
\]
\(\Pic^0(X)\) is divisible by a prime number \(p\) and \(T\) is a torsion group, so one can apply 2.1.3 to the restriction of the polynomial map (2.3.1) to \(\Pic^\tau(X)\). One thus obtains:


%%%%%%%%%%  scan idx 551  |  folio 538  %%%%%%%%%%
\SGAPrintPage{538}%

\sgaStableTarget{sga6.e10.prop.2-4-4}\textbf{Proposition 2.4.4.} \uline{The polynomial map} \(s_y\) (2.3.1) \uline{is constant on the cosets modulo} \(\Pic^\tau(X)\).

\emph{Proof}: If \(K\) is an algebraically closed extension of \(k\), one obtains a commutative diagram
\[
\begin{tikzpicture}[>=stealth,baseline=(current bounding box.center)]
\node (PicX) at (0,2) {\(\Pic(X)\)};
\node (Kk) at (3.7,2) {\(K_\bullet(\Spec(k))\)};
\node (Z) at (5.35,1) {\(\mathbb Z\)};
\node (Q) at (6.6,1) {\(\mathbb Q\)};
\node (PicXK) at (0,0) {\(\Pic(X_K)\)};
\node (KK) at (3.7,0) {\(K_\bullet(\Spec(K))\)};
\draw[->] (PicX) -- node[above] {\(s_y\)} (Kk);
\draw[->] (PicX) -- (PicXK);
\draw[->] (PicXK) -- node[above] {\(\underline{s}_y\)} (KK);
\draw[->] (Kk) -- node[sloped,above] {\(\sim\)} (Z);
\draw[->] (KK) -- node[sloped,above] {\(\sim\)} (Z);
\draw[right hook->] (Z) -- (Q);
\end{tikzpicture},
\]
where \(\underline{s}_y\) is the polynomial map on \(\Pic(X_K)\) corresponding to \(s_y\). Thus one may suppose that \(k\) is algebraically closed. The assertion then follows immediately from 2.1.3 and 2.4.3.

\sgaStableTarget{sga6.e10.rem.2-4-5}\textbf{Remark 2.4.5.} In Exp. XIII this result will be proved without using the representability of the functor \(\Pic_{X/k}\).

\sgaStableTarget{sga6.e10.sec.3}\section*{3. Projection formulas for the associated graded groups}

\sgaStableTarget{sga6.e10.sec.3.1}\subsection*{3.1.}
Consider a noetherian scheme \(X\). Denote by \(\Gro(X)\) the graded group associated with the \(\lambda\)-filtration of \(K^{\bullet}(X)\), and by \(\Grb(X)\) the graded group associated with the filtration 1.1.1 of \(K_\bullet(X)\). By 1.3.2, the composition law (1.3.1)
\[
 K^{\bullet}(X)\times K_\bullet(X)\longrightarrow K_\bullet(X)
\]
induces, on passing to the associated graded groups, composition laws
\sgaStableTarget{sga6.e10.eq.3-1-1}\begin{equation}\tag{3.1.1}
 \Gr^i(X)\times\Gr_j(X)\longrightarrow \Gr_{j-i}(X)\qquad(i,j\in\mathbb Z).
\end{equation}

